% =============================================================
%  Slab Allocator Crash Recovery  —  drop into your paper body
%  Requires: booktabs.  siunitx optional (for \SIrange).
% =============================================================

\subsection{Slab Allocator Crash Recovery}
\label{sec:slab-recovery}

Because the kernel slab allocator manages small-object pools
(\SIrange{32}{2048}{\byte}) whose metadata resides in CXL shared memory,
a machine crash in the middle of an allocation or deallocation can leave
the per-slab free list and its counter in an inconsistent state.
We address this with a lightweight \emph{per-CPU in-flight undo log}
stored in CXL shared memory, together with persistent pool metadata
that allows recovery to walk linked lists rather than scan all of CXL.

% ---------- data structure ----------
\subsubsection{Per-CPU Undo Log}

Each 2\,MB slab is fronted by a \texttt{slab\_header} that tracks a
linked free list (\texttt{free\_list\_head}) and a counter
(\texttt{current\_free\_cnt}).
Because each CPU can have at most one slab operation in flight
(protected by a per-order spinlock), we maintain a single 24-byte
\emph{undo log} per CPU in CXL shared memory
(\texttt{dsm\_meta->cxl\_slab\_meta[machine].cpu\_logs[cpu]}).
The log records:

\begin{itemize}
  \item \texttt{op} — the operation type
        (\textsc{None}, \textsc{Alloc}, or \textsc{Free});
  \item \texttt{slab\_addr} — pointer to the slab being modified;
  \item \texttt{old\_free\_head} — the value of \texttt{free\_list\_head}
        before the operation;
  \item \texttt{old\_free\_cnt} — the value of \texttt{current\_free\_cnt}
        before the operation.
\end{itemize}

\noindent
Crucially, the \texttt{slab\_header} itself is \emph{not} modified —
no per-slab log field is needed, preserving the original header size
and slot layout.  The per-CPU slab pool metadata
(\texttt{current\_slab}, \texttt{partial\_slab\_list},
\texttt{full\_slab\_list}) is also placed in CXL, so it survives
crashes and can be walked directly during recovery.

% ---------- protocol ----------
\subsubsection{Undo Protocol}

Both the allocation and deallocation paths follow the same three-step
protocol, executed entirely under the per-CPU spinlock:

\begin{enumerate}
  \item \textbf{Log.}\;
        Snapshot the slab address and current header fields into the
        per-CPU log, then set \texttt{op}.
        Persist with \texttt{clwb}\,+\,\texttt{sfence}.
  \item \textbf{Mutate.}\;
        Perform the actual slab mutation (advance the free-list pointer
        and update the counter; manage pool lists if the slab becomes
        full or returns from full).
        Persist the header with \texttt{clwb}\,+\,\texttt{sfence}.
  \item \textbf{Commit.}\;
        Set \texttt{op} back to \textsc{None} and persist
        (\texttt{clwb} only; the next operation's \texttt{sfence}
        orders it).
\end{enumerate}

\noindent
Because the commit step executes before the lock is released,
and the caller only receives the allocated pointer \emph{after} unlock,
a crash at any point during the critical section is safe to undo:
the caller could not have observed or used the result.

Figure~\ref{fig:slab-crash-semantics} summarizes the crash semantics.

\begin{figure}[t]
\centering
\small
\begin{tabular}{@{}llll@{}}
\toprule
\textbf{Crash point} & \textbf{Log state} & \textbf{Action} & \textbf{Outcome} \\
\midrule
Before log write       & \texttt{op=None}  & ---   & Op never started \\
After log, before mutate & \texttt{op=Alloc/Free} & Undo & Reverted; no effect \\
After mutate, before commit & \texttt{op=Alloc/Free} & Undo & Alloc: slot reclaimed; \\
                       &                   &       & Free: slot leaked \\
After commit           & \texttt{op=None}  & ---   & Op completed \\
\bottomrule
\end{tabular}
\caption{Crash semantics of the slab undo protocol.
         In all cases the slab is left in a consistent state
         (free-list length equals \texttt{current\_free\_cnt}).}
\label{fig:slab-crash-semantics}
\end{figure}

% ---------- recovery ----------
\subsubsection{Recovery Procedure}

On restart (or when a surviving machine re-attaches to the shared CXL
region), the recovery procedure executes three phases for each machine:

\paragraph{Phase 1: Undo in-flight operations.}
For each CPU, the procedure checks its per-CPU log in CXL.
If \texttt{op} is not \textsc{None}, the two header fields of the
referenced slab are restored from the log, \texttt{op} is cleared,
and the cache line is flushed.
This requires checking only $N_{\text{cpu}}$ log entries per machine —
no scan of the CXL page-frame metadata is needed.

\paragraph{Phase 2: Lock reset.}
All per-CPU, per-order slab spinlocks are re-initialized,
since no holder can still be alive after a crash.

\paragraph{Phase 3: Rebuild pool lists.}
Each per-CPU, per-order pool tracks three categories of slabs:
\texttt{current\_slab}, a \texttt{partial\_slab\_list}, and a
\texttt{full\_slab\_list} — all persisted in CXL.
After the undo phase, some slabs may have changed category
(e.g.\ a full slab whose in-flight allocation was undone is now partial).
The procedure collects all slabs from the three lists, then
re-classifies each:
fully free slabs are returned to the buddy allocator;
full slabs go to \texttt{full\_slab\_list};
and partial slabs are assigned as \texttt{current\_slab} or appended
to \texttt{partial\_slab\_list}.

% ---------- overhead ----------
\subsubsection{Overhead}

The protocol adds two \texttt{sfence} barriers per slab operation
on the hot path (one after the log write, one after the header
mutation).
The commit step elides its \texttt{sfence} — the next operation's
log barrier orders the clear.
When crash recovery is not needed, all instrumentation compiles to
empty macros via a build-time flag (\texttt{SLAB\_CRASH\_RECOVERY}),
incurring zero overhead.
The per-CPU log is fixed at 24\,bytes per CPU regardless of the number
of active slabs, and the \texttt{slab\_header} size is unchanged.

% ---------- correctness ----------
\subsubsection{Correctness Argument}

The undo log satisfies the following invariant:

\begin{quote}
\emph{At every point during execution, either
\textup{(a)}~\texttt{op\,=\,None} and the slab header is consistent, or
\textup{(b)}~\texttt{op\,$\neq$\,None} and the per-CPU log contains a
valid pre-image that, when restored, yields a consistent header.}
\end{quote}

\noindent
Consistency is defined as the free-list length equaling
\texttt{current\_free\_cnt}.
The invariant holds because:
(i)~the log is persisted \emph{before} any mutation (write ordering via
\texttt{sfence});
(ii)~the mutation is persisted \emph{before} the log is cleared;
and (iii)~only one operation can be in flight per CPU (enforced by the
spinlock).
We validate this with nine kernel-level unit tests that simulate
crashes at every intermediate state
(before mutation, after partial mutation, after full mutation)
for both allocation and deallocation across all seven size classes,
verifying post-recovery consistency on a two-machine DSM cluster.
